% =============================================================================
%  CHAPTER 18 -- PHASE 5: BASELINE MODEL COMPARISON
% =============================================================================
\section{Phase 5 --- Baseline model comparison: OLS, Ridge, Lasso, Elastic Net, PCR}
\label{ch:phase5}

\textbf{Reproduce it:} \code{python3 scripts/analysis/compare\_models.py \&\& python3
scripts/plotting/render\_model\_comparison.py}.\\
\textbf{Math used:} Chapter~\ref{ch:regularization} in full, plus
Sections~\ref{sec:ols:multicol} and \ref{sec:bt:walkforward}.\\
\textbf{Artifacts:} \code{scripts/model/baseline\_models.py},
\code{results/tables/model\_compare\_\{summary,sweep\}.csv},
\code{results/figures/model\_compare\_\{sharpe,ic,stability\}.png},
\code{report/tables/model\_compare.tex}.

\subsection{Goal}

The first genuine research question: \emph{does the choice of hedge estimator change the
traded result, net of cost?} Five estimators --- OLS, Ridge, Lasso, Elastic Net, and
principal-component regression --- through one identical causal, cost-aware pipeline, with
hyper-parameters chosen on a validation window and reported on a disjoint test window.

\subsection{Design: everything fixed except the estimator}

For each estimator the Phase~\ref{ch:phase4} hedge is reproduced exactly: at bar $t$, fit
$y=\alpha+x^{\top}\beta$ on the trailing window $W=120$ ending at $t-1$ (log prices),
standardise the bar-$t$ residual by the in-window residual $\sigma$ to get $z_t$, feed
$z_t$ through the two-band policy, execute one bar later, pay 10\,bps one-way. All five
therefore see the same data, the same costs, and the same execution delay. They differ
\emph{only} in how $\hat\beta$ is computed --- which is what makes the comparison an
experiment rather than a collection of anecdotes.

\paragraph{Why strategy-level metrics rather than $R^2$.} On log-price levels the in-sample
$R^2$ of any of these regressions is $\approx0.98$ and essentially uninformative
(Section~\ref{sec:prob:spurious}). So the comparison is made on:
\begin{itemize}
  \item \textbf{net Sharpe}, full sample (2010--2026) and test period (2022+);
  \item \textbf{out-of-sample information coefficient}: $\Corr(z_t,\ R^{\text{pos}}_{t,t+21})$,
    where $R^{\text{pos}}$ is the realized gross return of the position the signal implies
    --- a direct measure of predictive content in the units the strategy trades
    (Section~\ref{sec:nonlinear:target});
  \item \textbf{coefficient instability}: the time-average of each weight's
    $\sigma/|\mu|$ over the test period --- a proxy for how much the hedge churns, and
    therefore for how much turnover the estimator induces;
  \item \textbf{nonzero fraction} (did the penalty actually sparsify?) and
    \textbf{milliseconds per fit} (is the estimator affordable at scale?).
\end{itemize}

\paragraph{Hyper-parameter selection.} Validation window 2019--2021, objective = net
Sharpe, coarse grid: Ridge $\alpha\in\{0.1,1,10\}$; Lasso/ElasticNet
$\alpha\in\{10^{-4},10^{-3},10^{-2}\}$ (with \code{l1\_ratio}$=0.5$) on
column-standardised features; PCR $k\in\{1,2,3\}$. Selected: Ridge $\alpha=10$,
Lasso $\alpha=10^{-2}$, ElasticNet $\alpha=10^{-2}$, PCR $k=1$. All reported test figures
are then strictly out-of-sample on 2022--2026. The full sweep is kept in
\code{results/tables/model\_compare\_sweep.csv} so the selection is auditable --- and so
that every grid point can be counted as a hypothesis in Phase~\ref{ch:phase10}.

\begin{remark}[Honesty: the validation window was a bad year to select on]
Every selected configuration has a \emph{negative} validation net Sharpe (OLS $-0.86$,
Ridge $-0.63$, Lasso $-0.38$, ElasticNet $-0.73$, PCR $-0.11$): the 2019--2021 window is
one in which the strategy loses money at 10\,bps for every estimator. Selecting the
``least bad'' on a uniformly negative window is a weak selection rule, and the table shows
it: the validation ranking (PCR best, Lasso second) does not match the test ranking (Ridge
best). Reporting the sweep rather than only the selected row is what makes this visible.
\end{remark}

\subsection{Results}

\input{tables/model_compare.tex}

\begin{figure}[H]\centering
\includegraphics[width=0.92\linewidth]{figures/model_compare_sharpe.png}
\caption{Gross versus net annualised Sharpe by hedge estimator, out-of-sample on
2022--2026 (hyper-parameters chosen on 2019--2021) with the full-sample values for
context.}
\label{fig:mc_sharpe}
\end{figure}

\begin{figure}[H]\centering
\includegraphics[width=0.92\linewidth]{figures/model_compare_ic.png}
\caption{Out-of-sample information coefficient: correlation of the traded $z_t$ with the
realized 21-bar gross position return, test period 2022--2026. Every estimator is
positive; the magnitudes are $0.125$--$0.211$.}
\label{fig:mc_ic}
\end{figure}

\begin{figure}[H]\centering
\includegraphics[width=0.95\linewidth]{figures/model_compare_stability.png}
\caption{Causal walk-forward hedge weights over the full sample: the weight on
\code{CRM} and the time-varying hedge norm $\lVert\beta_t\rVert_2$ for each estimator.
The OLS trace is visibly the noisiest; ridge and PCR are smooth.}
\label{fig:mc_stab}
\end{figure}

Table~\ref{tab:model_compare} is the generated full record; the curated extract below
is the same data read column by column (all net figures at 10\,bps one-way):

\begin{center}
\renewcommand{\arraystretch}{1.2}
\small
\begin{tabular}{lrrrrrrrr}
\toprule
Estimator & \multicolumn{2}{c}{Full sample} & \multicolumn{2}{c}{Test 2022+} & OOS IC & Instab. & ms/fit & Trades \\
& gross SR & net SR & gross SR & net SR & ($h{=}21$) & & & \\
\midrule
OLS (none)          & $0.391$ & $\mathbf{+0.050}$ & $-0.017$ & $-0.302$ & $0.125$ & $0.998$ & $0.67$ & 48 \\
Ridge ($\alpha{=}10$) & $0.341$ & $\mathbf{+0.094}$ & $+0.198$ & $-0.031$ & $0.150$ & $0.935$ & $1.50$ & 45 \\
Lasso ($\alpha{=}10^{-2}$) & $0.149$ & $+0.007$ & $+0.060$ & $-0.103$ & $0.145$ & $1.098$ & $1.21$ & 44 \\
ElasticNet          & $0.163$ & $-0.032$ & $+0.038$ & $-0.145$ & $0.157$ & $1.038$ & $1.20$ & 46 \\
PCR ($k{=}1$)        & $0.166$ & $+0.026$ & $-0.146$ & $-0.250$ & $\mathbf{0.211}$ & $\mathbf{0.851}$ & $1.90$ & 34 \\
\bottomrule
\end{tabular}
\end{center}
(Trades are the test-period round trips; the full-sample counts are in
Table~\ref{tab:model_compare}.)

\subsection{Findings}

\begin{enumerate}[label=\textbf{M\arabic*.}]
  \item \textbf{The Python OLS walk-forward reproduces the C++ baseline.} Full-sample net
    Sharpe $\approx0.050$ matches the Phase~\ref{ch:phase4} engine to three decimals.
    Two implementations, same experiment --- so every difference below is a model effect.
  \item \textbf{Regularization improves \emph{net} performance by reducing turnover, not
    by finding signal.} Ridge ($\alpha{=}10$) lifts the full-sample net Sharpe from
    $0.050$ to $0.094$ while its \emph{gross} Sharpe is slightly \emph{lower} than OLS
    ($0.341$ vs $0.391$). The mechanism is visible in the coefficient-instability column
    ($0.998\to0.935$) and in Figure~\ref{fig:mc_stab}: the L2 shrinkage damps the
    small-eigenvalue directions of $X^{\top}X$ \eqref{eq:ridge:spectral}, which are
    exactly the directions in which the OLS hedge oscillates from window to window. A
    smoother hedge means fewer position changes, less traded notional, less cost --- the
    arithmetic of \eqref{eq:costdrag}. On the test period Ridge is also the
    least-negative net Sharpe ($-0.031$ vs OLS $-0.302$).
  \item \textbf{Aggressive sparsification hurts a hedging regression.} Lasso and
    ElasticNet, driven to heavy penalties by the validation rule, degrade net performance
    ($+0.007$ and $-0.032$ full sample). This is the empirical confirmation of the
    inductive-bias argument of Section~\ref{sec:reg:comparison}: in a three-name
    software basket every name carries shared factor exposure, so setting a weight to
    exactly zero removes a hedge leg the spread needs. Note that the nonzero fraction is
    $1.0$ for both --- at the selected $\alpha$ the penalty shrinks without fully
    sparsifying, and the damage comes from the shrinkage being applied in the wrong
    geometry (L1 treats each name independently) rather than from dropping names outright.
  \item \textbf{PCA has the strongest raw signal and the weakest realized hedge.} The
    one-component PCR residual gives the highest OOS IC ($0.211$ vs $0.125$ for OLS) and
    the lowest coefficient instability ($0.851$), yet its full-sample net Sharpe ($0.026$)
    is below OLS and its test-period net Sharpe ($-0.250$) is the worst but one. The
    explanation is the eigenstructure argument of Section~\ref{sec:pca:derivation}: the
    first principal component is the common sector/market mode, so projecting onto it
    removes the market drift cleanly (hence a crisp, informative signal) but discards the
    idiosyncratic relative moves that the mean-reversion trade actually monetises, and
    under-hedges the names' common drift --- leaving the book exposed to a residual factor
    it no longer measures. High IC, poor transfer coefficient
    (Section~\ref{sec:nonlinear:fundamentallaw}).
  \item \textbf{The predictive content is real but small, and the test period is adverse.}
    Every estimator has a positive OOS IC on 2022--2026 ($0.125$--$0.211$), i.e.\ the
    standardised spread genuinely predicts the \emph{direction} of convergence out of
    sample. Yet every estimator is net-negative there, because 2022 was the worst regime
    for this book (Phase~\ref{ch:phase7}: net Sharpe $-1.31$ in the 2022 selloff window):
    rising rates and the growth-stock rotation broke the software/semis relationships the
    baskets are built on, and 10\,bps per side turned a small edge into a loss.
\end{enumerate}

\subsection{Computational cost}

Measured fit times on a $120\times3$ window in sklearn: $0.67$\,ms for OLS,
$1.20$--$1.90$\,ms for the regularised variants --- a factor of $\approx2$--$3$, entirely
immaterial to model choice at these sizes (the deciding factors are statistical). For
scale comparison, the dependency-free C++ path performs the equivalent update-and-solve in
$\approx0.6\,\mu$s (Table~\ref{tab:engine_latency}), i.e.\ $\approx1{,}000\times$ faster,
which is what makes the $10^{5}$--$10^{6}$-backtest universe scan recommended by
Chapter~\ref{ch:multipletesting} feasible at all.

\subsection{Honest reading}

The finding ``ridge beats OLS net of cost'' is directionally robust and mechanistically
explained, and it is the first result in this project that is a genuine engineering
conclusion rather than a market claim: \emph{shrink the hedge, trade less, pay less}.
Its magnitude is not established: the DM test of Phase~\ref{ch:phase10} between the
Ridge and OLS hedges gives a mean daily difference of $0.32$\,bps with a HAC standard
error of $1.14$\,bps ($p=0.78$) --- economically meaningful, statistically invisible.
And the ``best'' estimator depends on which window you look at: Ridge is best full-sample,
PCR has the best IC, OLS-w60 is best on 2022+ (Phase~\ref{ch:phase8}). Selecting the
winner after seeing all three is exactly the snooping that $M=44$ accounts for.

\subsection{Exit gate --- status}

\begin{itemize}
  \item[\checkmark] all five linear variants run through a common pipeline;
  \item[\checkmark] OOS metrics and stability time series plotted;
  \item[\checkmark] hyper-parameter sweep saved and auditable;
  \item[\checkmark] report section written from the generated tables, with the
    honest reading above included rather than only the best row.
\end{itemize}
